\documentclass[11pt,a4paper]{article}

\usepackage[T1]{fontenc}
\usepackage[utf8]{inputenc}
\usepackage[margin=2.5cm]{geometry}
\usepackage{booktabs}
\usepackage{tabularx}
\usepackage{enumitem}
\usepackage{xcolor}
\usepackage{listings}
\usepackage[hidelinks]{hyperref}

\setlength{\emergencystretch}{2em}

\lstset{
  basicstyle=\ttfamily\small,
  breaklines=true,
  frame=single,
  backgroundcolor=\color{gray!8},
  columns=fullflexible
}

\title{SOC Home Lab\\[0.3em]
  \large Cyberattack Detection, Monitoring and Incident Response}
\author{ABDELthe1}
\date{November 2026}

\begin{document}

\maketitle

\begin{abstract}
An enterprise-inspired Security Operations Center (SOC) built at home: a Wazuh
SIEM and Suricata IDS monitor a small Active Directory environment that is
attacked in a controlled way from a separate Parrot Security workstation. The
full cycle is practiced and documented: attack, collect, detect, alert,
investigate, contain, remediate, report. This report covers Phase~1 (network
preparation), Phase~2 (VM deployment), Phase~3 (Active Directory), Phase~4
(Wazuh SIEM and agents), Phase~5 (endpoint monitoring) and Phase~6 (Suricata
IDS), which are complete. Phase~7 (attack simulations) is next.
\end{abstract}

\tableofcontents
\newpage

%--------------------------------------------------------------------
\section{Objectives}

\begin{enumerate}
  \item Generate controlled cyberattacks against an isolated lab.
  \item Collect endpoint and network telemetry centrally.
  \item Detect suspicious activity with default and custom rules.
  \item Investigate alerts: source, target, user, timeline, impact.
  \item Contain and remediate (NIST incident response lifecycle).
  \item Document every incident with evidence and a professional report.
\end{enumerate}

%--------------------------------------------------------------------
\section{Architecture}

\begin{table}[h]
\centering
\begin{tabularx}{\textwidth}{lX}
\toprule
Component & Role \\
\midrule
PC1 (Windows 11 + VMware Workstation) & Hosts 4 VMs \\
PC2 (Parrot Security 7.4) & Red Team / attacker workstation (physical machine) \\
VM1 SOC Server & Wazuh Manager + Indexer + Dashboard, Suricata IDS \\
VM2 Windows Server 2019 & Domain Controller (AD DS, DNS), Sysmon, Wazuh agent \\
VM3 Windows 10 Pro & Domain-joined employee workstation, Sysmon, Wazuh agent \\
VM4 Ubuntu Server 24.04 & Linux server, SSH, auditd, UFW, Wazuh agent \\
\bottomrule
\end{tabularx}
\caption{Lab components}
\end{table}

\paragraph{Architecture change.} The Domain Controller uses Windows Server 2019
instead of 2022, because the 2019 ISO was available, it is lighter, and it is
widely used in AD attack and detection labs.

%--------------------------------------------------------------------
\section{Network design}

Neither PC has a second NIC, so the lab subnet is created by \textbf{IP
aliasing} on the existing interface. The lab is logically separated, not
physically separated.

\begin{table}[h]
\centering
\begin{tabular}{ll}
\toprule
Item & Value \\
\midrule
Home LAN (unchanged) & 192.168.1.0/24, gateway 192.168.1.1 \\
Lab subnet & 10.10.10.0/24, \textbf{no default gateway} \\
VMware network & VMnet0, bridged to the physical Ethernet adapter \\
Internet for VMs & Temporary NAT adapter (VMnet8), updates only \\
\bottomrule
\end{tabular}
\caption{Network parameters}
\end{table}

\begin{table}[h]
\centering
\begin{tabular}{lll}
\toprule
Device & Home LAN & Lab IP \\
\midrule
PC1 host & 192.168.1.8 (static) & 10.10.10.1 (alias) \\
PC2 Parrot & 192.168.1.12 (DHCP) & 10.10.10.50 (alias) \\
Windows Server 2019 (DC) & -- & 10.10.10.10 \\
SOC Server & -- & 10.10.10.20 \\
Windows Client & -- & 10.10.10.30 \\
Ubuntu Server & -- & 10.10.10.40 \\
\bottomrule
\end{tabular}
\caption{IP addressing}
\end{table}

\paragraph{Isolation rules.}
\begin{itemize}
  \item Lab IPs have no gateway, so the router never routes 10.10.10.0/24.
  \item Attack tools are scoped to 10.10.10.0/24 only.
  \item No home device is placed in 10.10.10.0/24.
  \item Nothing in the lab is exposed to the internet.
\end{itemize}

\paragraph{Known limitation.} The lab shares the physical wire with home
devices. A direct cable between the PCs or a dedicated switch is a future
improvement.

\paragraph{IDS visibility.} Suricata runs on a dedicated capture vNIC of the SOC Server
(\texttt{ens38}, promiscuous mode, no IP) attached to VMnet0, so it sees traffic that reaches
PC1's NIC and traffic on the virtual switch. Visibility is validated before any
detection is trusted.

%--------------------------------------------------------------------
\section{VM inventory}

\begin{table}[h]
\centering
\small
\begin{tabularx}{\textwidth}{lXcccl}
\toprule
VM & OS & vCPU & RAM & Disk & Lab IP \\
\midrule
SOC-Server & Ubuntu Server 24.04.3 LTS & 4 & 8 GB & 80 GB (root LV 78 GB) & 10.10.10.20 \\
DC01 & Windows Server 2019 Standard Evaluation (Desktop Experience) & 2 & 5 GB & 60 GB & 10.10.10.10 \\
WIN-CLIENT & Windows 10 Pro & 2 & 5 GB & 60 GB & 10.10.10.30 \\
UBUNTU-SRV & Ubuntu Server 24.04.3 LTS & 2 & 2 GB & 30 GB (root LV 28 GB) & 10.10.10.40 \\
\bottomrule
\end{tabularx}
\caption{Virtual machines (domain: \texttt{SOC-LAB.LOCAL})}
\end{table}

Domain \texttt{SOC-LAB.LOCAL}: AD DS and DNS on DC01, DNS forwarder 8.8.8.8.
WIN-CLIENT is joined to the domain. The Domain Controller (DC01) runs Windows
Server 2019.

Each VM has a lab adapter on VMnet0 (static lab IP, no gateway) and a
temporary NAT adapter (VMnet8) for updates only. All four VMs reach each other
on 10.10.10.0/24.

%--------------------------------------------------------------------
\section{Phase 1: Network preparation}

\subsection{Done}
\begin{itemize}
  \item PC1: static 192.168.1.8/24 (gateway and DNS 192.168.1.1) plus alias
        10.10.10.1/24 on \texttt{Ethernet}.
  \item PC1: firewall rule allowing inbound ICMPv4 from 10.10.10.0/24.
  \item PC2: alias 10.10.10.50/24 on its wired interface, persisted in the
        NetworkManager profile.
  \item Cross-ping verified in both directions; Parrot kept its internet access.
  \item VMware: VMnet0 added in the Virtual Network Editor; VMware Bridge
        Protocol enabled on the physical adapter.
\end{itemize}

\subsection{Problems and lessons learned}
\begin{enumerate}
  \item Adding the alias with \texttt{New-NetIPAddress} on a DHCP-enabled
        adapter disabled DHCP and removed the 192.168.1.8 lease, so PC1 lost
        internet. Fix: set a static primary address with \texttt{netsh}, then
        add the alias. Windows cannot keep a DHCP primary and a static
        secondary address on one adapter.
  \item \texttt{ipconfig} did not list the new alias immediately because
        Windows hides tentative addresses until duplicate-address detection
        finishes. \texttt{Get-NetIPAddress} showed the real state.
  \item On Parrot, \texttt{nmcli con up} would have dropped the SSH session.
        The alias was added live with \texttt{ip addr add} first, and the
        profile was changed without reactivating it.
\end{enumerate}

%--------------------------------------------------------------------
\section{Phase 2: VM deployment}

\subsection{SOC-Server}
\begin{itemize}
  \item Ubuntu Server 24.04.3 LTS, 4 vCPU, 8 GB RAM, 80 GB disk.
  \item \texttt{ens33} = VMnet0, static 10.10.10.20/24;
        \texttt{ens34} = NAT (DHCP) for updates.
  \item Root LV extended from 39 GB to 78 GB.
  \item SSH works from Parrot and PC1. Fully updated.
  \item Snapshot \texttt{01-clean-install-network-ok} taken.
  \item Wazuh installed in Phase~4; Suricata installed in Phase~6.
\end{itemize}

\subsection{DC01}
\begin{itemize}
  \item Windows Server 2019 Standard Evaluation (Desktop Experience),
        2 vCPU, 5 GB RAM, 60 GB disk.
  \item LAB adapter (VMnet0) static 10.10.10.10/24, no gateway; NAT adapter
        for updates.
  \item Pings to 10.10.10.1, 10.10.10.20 and 10.10.10.50 succeed.
  \item Promoted to Domain Controller in Phase~3.
\end{itemize}

\subsection{WIN-CLIENT}
\begin{itemize}
  \item Windows 10 Pro, 2 vCPU, 5 GB RAM, 60 GB disk.
  \item LAB adapter (VMnet0) static 10.10.10.30/24, no gateway, DNS
        10.10.10.10; second adapter on NAT, not registered in DNS.
  \item Joined to \texttt{SOC-LAB.LOCAL}; computer object in
        \texttt{OU=Computers,\allowbreak OU=SOC-LAB,\allowbreak DC=SOC-LAB,\allowbreak DC=LOCAL}.
  \item Login as \texttt{SOC-LAB\textbackslash IT\_User} verified
        (\texttt{whoami} and domain name checked).
\end{itemize}

\subsection{UBUNTU-SRV}
\begin{itemize}
  \item Ubuntu Server 24.04.3 LTS, 2 vCPU, 2 GB RAM, 30 GB disk; root LV
        extended to 28 GB.
  \item \texttt{ens33} = VMnet0, static 10.10.10.40/24 (no gateway);
        \texttt{ens34} = NAT (DHCP).
  \item OpenSSH installed; SSH works from Parrot.
  \item Ping verified to 10.10.10.10, 10.10.10.20 and 10.10.10.50.
\end{itemize}

All four VMs are created and reach each other on 10.10.10.0/24.

\subsection{Problems and lessons learned}
\begin{enumerate}
  \item \textbf{ICMP redirects and packet loss.} Duplicate ping replies
        (TTL 64 and 63), ICMP redirects and 50\% packet loss between the VM
        and Parrot. Cause: IP forwarding was enabled on PC1's
        \texttt{Ethernet} interface. Fix:
\begin{lstlisting}
Set-NetIPInterface -InterfaceAlias "Ethernet" -AddressFamily IPv4 -Forwarding Disabled
sudo ip route flush cache   # on the Linux machines
\end{lstlisting}
        After the fix: no redirects, no duplicates, 0\% loss, TTL 64.
  \item \textbf{Root LV size.} The Ubuntu installer uses only about half of
        the disk for the root LV by default. Fix:
\begin{lstlisting}
sudo lvextend -r -l +100%FREE /dev/mapper/ubuntu--vg-ubuntu--lv
\end{lstlisting}
  \item \textbf{Bridged (Automatic).} A VM network adapter set to
        ``Bridged (Automatic)'' can bind to the wrong host adapter. Use
        Custom: VMnet0 instead.
  \item \textbf{Windows edition.} Easy Install in the VMware wizard
        preselected ``Datacenter Core'' (no GUI). Pick the Desktop Experience
        edition.
  \item \textbf{One-off TTL 128.} TTL 128 appeared once on a ping from DC01
        to the SOC-Server. Later checks showed ttl=64, a correct ARP table and
        forwarding disabled, so it was treated as a one-off.
  \item \textbf{Missing lab adapter.} WIN-CLIENT was installed with only a
        NAT adapter. The lab adapter on VMnet0 was hot-added afterwards in VM
        Settings.
  \item \textbf{Installer SSH host key.} SSH to 10.10.10.40 during the Ubuntu
        install reached the installer's temporary SSH server (different host
        key and password), causing ``Permission denied'' and a changed host key
        warning after the reboot. Fix: \texttt{ssh-keygen -R} for the old
        entry. Lesson: do not SSH into a VM until the install is finished and
        rebooted.
  \item \textbf{Keyboard layout.} On Linux, a password typed during the
        install used the US keyboard layout. If SSH rejects it, reset it from
        the VM console with \texttt{sudo passwd}.
\end{enumerate}

%--------------------------------------------------------------------
\section{Phase 3: Active Directory}

\subsection{Done}
\begin{itemize}
  \item DC01 (Windows Server 2019 Standard Evaluation, Desktop Experience)
        runs AD DS and DNS for the domain \texttt{SOC-LAB.LOCAL}, with DNS
        forwarder 8.8.8.8.
  \item OU structure: \texttt{SOC-LAB} containing Users, Administrators,
        IT Department, HR Department, Finance Department, Computers, Servers.
  \item Users created and enabled: \texttt{IT\_User} (IT Department),
        \texttt{HR\_User} (HR Department), \texttt{Finance\_User} (Finance
        Department), \texttt{Standard\_User} (Users). Built-in Administrator
        exists.
\end{itemize}

\subsection{Problems and lessons learned}
\begin{enumerate}
  \item \textbf{NAT adapter registered in DNS.} DC01 registered its NAT
        adapter address in DNS, so \texttt{soc-lab.local} resolved to two IPs.
        Fix:
\begin{lstlisting}
Set-DnsClient -InterfaceAlias NAT -RegisterThisConnectionsAddress $false
dnscmd /ResetListenAddresses 10.10.10.10
\end{lstlisting}
        then delete the stale A records in \texttt{SOC-LAB.LOCAL} and
        \texttt{\_msdcs.SOC-LAB.LOCAL}, restart DNS and Netlogon, and run
        \texttt{nltest /dsregdns}. The same setting was applied on WIN-CLIENT
        before the domain join. Lesson: multi-homed lab machines must not
        register the temporary NAT adapter in AD DNS.
\end{enumerate}

%--------------------------------------------------------------------
\section{Phase 4: Wazuh SIEM and agents}

\subsection{Done}
\begin{itemize}
  \item Wazuh 4.14.10 all-in-one installed on SOC-Server (10.10.10.20).
  \item Agents registered on DC01, WIN-CLIENT and UBUNTU-SRV.
  \item VMs moved from the USB disk to NVMe
        (\texttt{C:\textbackslash VMs\textbackslash project}) to fix 100\%
        disk activity and slowness.
\end{itemize}

\subsection{Problems and lessons learned}
\begin{enumerate}
  \item \textbf{Dashboard API errors.} The dashboard showed ``No API
        available'' / HTTP 429. Cause: retries during start-up. Fix: wait
        until the services are fully up and use one browser tab.
\end{enumerate}

%--------------------------------------------------------------------
\section{Phase 5: Endpoint monitoring (Sysmon, auditd, FIM)}

\subsection{Windows (WIN-CLIENT, DC01)}
\begin{itemize}
  \item Sysmon (SwiftOnSecurity config), collected by the Wazuh agent via an
        \texttt{eventchannel} localfile.
  \item Windows audit policies set with \texttt{auditpol}: Logon, Account
        Lockout, Credential Validation, User/Group Management, Process
        Creation; Kerberos on DC01.
  \item Validated: a failed logon to DC01 from WIN-CLIENT produced Event 4625
        and the rule ``Logon Failure - Unknown user or bad password''
        (level~5).
\end{itemize}

\subsection{Linux (UBUNTU-SRV)}
\begin{itemize}
  \item auditd rules (\texttt{Configuration/\allowbreak auditd/\allowbreak soc.rules},
        listed below) and FIM in realtime on
        \texttt{/home/\allowbreak socadmin/\allowbreak fim-test}.
  \item Validated: a file added in the FIM folder raised rule 554;
        \texttt{useradd} / \texttt{userdel} produced rules 80781, 80791 and
        80792.
\end{itemize}

Final auditd rules:
\begin{lstlisting}
-w /etc/passwd -p wa -k audit-wazuh-w
-w /etc/shadow -p wa -k audit-wazuh-w
-w /etc/sudoers -p wa -k audit-wazuh-w
-w /etc/sudoers.d/ -p wa -k audit-wazuh-w
-a always,exit -F arch=b64 -F euid=0 -F auid>=1000 -F auid!=4294967295 -S execve -k audit-wazuh-c
\end{lstlisting}

\subsection{Problems and lessons learned}
\begin{enumerate}
  \item \textbf{Noisy root \texttt{execve} rule.} The initial auditd rule
        logging every root \texttt{execve} produced about 1,600 events in
        15 minutes. Fix: restricted to \texttt{auid>=1000} and switched the
        watch keys to \texttt{audit-wazuh-w} so the Wazuh built-in rules
        fire.
  \item \textbf{Time zones.} VM time zones differ (DC01 is on Pacific time).
        UTC timestamps were verified correct.
\end{enumerate}

%--------------------------------------------------------------------
\section{Phase 6: Suricata IDS}

\subsection{Done}
\begin{itemize}
  \item Suricata 8.0.8 (OISF PPA) installed on SOC-Server.
  \item Third vNIC \texttt{ens38} on VMnet0 for capture: no IP, promiscuous,
        GRO/LRO/TSO/GSO off, made persistent with the systemd unit
        \texttt{capture-nic.service}.
  \item \texttt{HOME\_NET} set to the four victim hosts only (10.10.10.10,
        .20, .30, .40), so the Parrot attacker 10.10.10.50 counts as
        \texttt{EXTERNAL\_NET}.
  \item BPF filter \texttt{net 10.10.10.0/24} on the capture interface.
  \item \texttt{eve.json} forwarded to Wazuh via a manager localfile
        (\texttt{log\_format json}). Wazuh matches all Suricata events with
        built-in rule 86601.
\end{itemize}

\subsection{Rule sources}
\begin{itemize}
  \item Emerging Threats Open via \texttt{suricata-update}: 69,559 rules
        loaded, 53,582 enabled, merged into
        \texttt{/var/lib/\allowbreak suricata/\allowbreak rules/\allowbreak suricata.rules}.
  \item Suricata built-in rules from
        \texttt{/usr/share/\allowbreak suricata/\allowbreak rules} (decoder,
        app-layer and stream events).
  \item No custom rules yet (Phase~8).
\end{itemize}

\subsection{Validation}
\texttt{nmap -sS -sV -O} from Parrot to UBUNTU-SRV produced the following
alerts, visible in Suricata (\texttt{fast.log} / \texttt{eve.json}) and in
Wazuh (rule 86601, groups \texttt{ids}/\texttt{suricata},
\texttt{in\_iface} \texttt{ens38}):

\begin{table}[h]
\centering
\begin{tabular}{rl}
\toprule
SID & Signature \\
\midrule
2001219 & ET SCAN Potential SSH Scan \\
2018489 & ET SCAN NMAP OS Detection Probe \\
2200025 & SURICATA ICMPv4 unknown code \\
2200130 & SURICATA TCP wscale over maximum \\
\bottomrule
\end{tabular}
\caption{Suricata alerts from the Nmap scan}
\end{table}

\subsection{Problems and lessons learned}
\begin{enumerate}
  \item \textbf{Host internet traffic captured.} Suricata also captured the
        host's internet traffic through the bridged segment (privacy and
        noise). Fix: BPF filter \texttt{net 10.10.10.0/24}.
  \item \textbf{Manager restart load.} A Wazuh manager restart caused a load
        spike (load average $\sim$27 on 4 vCPUs) and a dashboard API timeout
        for a few minutes.
  \item \textbf{Manager events.} Manager events (agent 000) are only visible
        in the global Threat Hunting view.
  \item \textbf{Noisy stream alerts.} Generic SURICATA STREAM alerts are
        noisy; they will be tuned in Phase~8.
\end{enumerate}

%--------------------------------------------------------------------
\section{Phase 7: Attack simulations (in progress)}

\subsection{Scenario 4: AD enumeration and password spray (complete)}

AD reconnaissance and password spraying with NetExec (\texttt{nxc}) from
Parrot (10.10.10.50) against DC01 (10.10.10.10), domain
\texttt{SOC-LAB.LOCAL}, on 2026-11-24 at about 14:35--14:36. Lab
10.10.10.0/24 only, authorized.

\paragraph{Reconnaissance.}
\begin{itemize}
  \item \texttt{nxc smb} enumeration confirmed DC01: Windows 10/Server 2019
        build 17763, SMB signing enabled, Null Auth allowed on SMB.
  \item Anonymous SMB user enumeration returned no users.
  \item Anonymous LDAP bind was correctly rejected by AD
        (``operationsError, successful bind must be completed'').
\end{itemize}

\paragraph{Password spray.} 5 usernames x 4 common passwords = 20 attempts
via SMB/NTLM with \texttt{-{}-continue-on-success}. 0 valid credentials
found.

\paragraph{Detections observed.}
\begin{table}[h]
\centering
\begin{tabular}{rrlr}
\toprule
Rule & Level & Description & Hits \\
\midrule
60122 & 5 & Logon Failure - Unknown user or bad password & 18 \\
60204 & 10 & Multiple Windows Logon Failures & 2 \\
\bottomrule
\end{tabular}
\caption{Wazuh alerts from the password spray}
\end{table}

\begin{itemize}
  \item Rule 60122 events: source IP 10.10.10.50, authentication package
        NTLM, logon type 3 (network).
  \item Rule 60204 confirms that Wazuh's built-in frequency correlation
        detected the spray pattern (same mechanism as the SSH brute-force
        scenario).
  \item No account lockout (event ID 4740): the attempt volume stayed under
        the default lockout threshold, consistent with realistic low-and-slow
        spray behaviour.
\end{itemize}

\paragraph{Phase 8 finding: Kerberos gap.} The default Wazuh ruleset
(\texttt{0580-win-\allowbreak security\_\allowbreak rules.xml},
\texttt{0220-msauth\_\allowbreak rules.xml})
covers event ID 4769 (TGS, rules 60106 / 60131) but has no Kerberos-specific
rule for 4768 (TGT requested) or 4771 (Kerberos pre-authentication failed).
These events were caught generically by rule 60104 (Windows audit failure,
level~5), but without Kerberos-specific identification or spray correlation.
A Kerberos-native password spray would therefore surface only as
undifferentiated level~5 audit failures, with no spray escalation; detection
in this run relied entirely on the NTLM/SMB path (rules 60122 / 60204). The
Phase~8 fix now adds this. Fix implemented and validated: custom local rules
100005 (4771), 100006 (4768) and 100007 (frequency correlation, level~12);
a live kerbrute password spray produced 5 x level~8 rule 100005 and
1 x level~12 rule 100007 (see \texttt{phase8-detection-engineering.md}).

\paragraph{MITRE ATT\&CK.} T1087.002 Domain Account Discovery, T1110.003
Password Spraying, T1021.002 SMB/Windows Admin Shares (attempted access
vector).

\paragraph{Phase 8 fixes from the Scenario 5 plan.}
\begin{itemize}
  \item \textbf{SUID detection (\texttt{chmod}): fix implemented and
        validated.} FIM extended to \texttt{/usr/local/bin},
        \texttt{/usr/bin} and \texttt{/usr/sbin}; rule 100009 (level~12)
        fired on \texttt{chmod u+s} (\texttt{rwxr-xr-x} $\rightarrow$
        \texttt{rwsr-xr-x}) and did not fire on \texttt{chmod u-s}.
        Caveat: only the \texttt{rws} pattern is matched, so SUID on a
        non-executable file (\texttt{rwS}) is not caught.
  \item \textbf{History tampering detection: fix implemented and
        validated.} FIM extended to \texttt{.bash\_history}; rule 100010
        (level~12) fired on modification (rule 550) and deletion (rule 553).
        Caveat: it fires on any \texttt{.bash\_history} change, not only
        malicious clearing, and the in-memory \texttt{history -c} and
        \texttt{HISTFILE} tampering remain undetected.
  \item \textbf{Privilege group modification (\texttt{usermod} /
        \texttt{gpasswd}): fix implemented and validated.} Rule 100008
        (level~10) on execve audit events fired on a live \texttt{sudo
        usermod} group modification, identifying the command and acting uid.
        Caveat: it escalates on any \texttt{usermod} / \texttt{gpasswd} run,
        not only ones targeting a privileged group; the target group must be
        checked in the raw log.
\end{itemize}
Details: \texttt{phase8-detection-engineering.md}.

%--------------------------------------------------------------------
\section{Open items (Phase 7 and later)}

\begin{itemize}
  \item Phase~7: remaining attack scenarios.
  \item Phase~8: custom detection rules and tuning of the SURICATA STREAM
        alerts.
\end{itemize}

%--------------------------------------------------------------------
\section{Planned attack scenarios}

All attacks stay inside the authorized lab network 10.10.10.0/24.

\begin{table}[h]
\centering
\small
\newcolumntype{L}{>{\raggedright\arraybackslash}X}
\begin{tabularx}{\textwidth}{cLLLL}
\toprule
\# & Scenario & Source $\rightarrow$ Target & Tools & Detection \\
\midrule
1 & Brute force & Parrot $\rightarrow$ Ubuntu / DC & Hydra & Auth logs + Wazuh rules \\
2 & Port scanning & Parrot $\rightarrow$ DC / Ubuntu & Nmap & Suricata \\
3 & Suspicious PowerShell & Windows client & PowerShell & Sysmon + Wazuh \\
4 & AD attacks & Parrot $\rightarrow$ DC & Impacket, NetExec, BloodHound & Windows Security + Sysmon + Wazuh \\
5 & Linux privilege indicators & Parrot / local $\rightarrow$ Ubuntu & CLI, SSH & auditd + Wazuh \\
6 & Incident response & Alerts from 1--5 & Wazuh dashboard & Full investigation and report \\
\bottomrule
\end{tabularx}
\caption{Attack scenarios}
\end{table}

A detection counts as successful only when the event is generated, collected,
analyzed and verified.

%--------------------------------------------------------------------
\section{Project status}

\begin{table}[h]
\centering
\begin{tabularx}{\textwidth}{clX}
\toprule
Phase & Description & Status \\
\midrule
1 & Infrastructure and network preparation & Done \\
2 & VM deployment & Done \\
3 & Active Directory & Done \\
4 & Wazuh SIEM and agents & Done \\
5 & Endpoint monitoring (Sysmon, auditd, FIM) & Done \\
6 & Suricata IDS & Done \\
7 & Attack simulations & Next \\
8 & Detection engineering & Pending \\
9 & Incident response & Pending \\
10 & Final documentation & Pending \\
\bottomrule
\end{tabularx}
\caption{Phase status}
\end{table}

%--------------------------------------------------------------------
\section{Scope and ethics}

All attacks target only the authorized lab network 10.10.10.0/24, owned by the
author. No lab system is exposed to the internet.

\end{document}
